% 3-1-intro.tex
%  Budget: 1pp.  Chapter roadmap.
%  Source map: no external claims; this section only announces the chapter's
%   own structure. The single substantive commitment (fair-by-construction
%   comparison under a published protocol) -> lago_forecasting_2021.
%  Numerals: Persian digits in body text; \lr{} only for Latin identifiers
%   and for em-dashes. See NUMERAL_CONVENTION (supersedes HANDOFF 1.1).

\section{مقدمه}
\label{sec:3-1}

این فصل روشِ پژوهش را چنان شرح می‌دهد که خواننده بتواند آن را بازتولید کند. ترتیب
بخش‌ها از داده به سوی مدل و سپس به سوی سنجش پیش می‌رود، زیرا هر حلقه از این زنجیره
اعتبار حلقه‌های پس از خود را تعیین می‌کند: ماتریس ویژگیِ آلوده، هر مدلی را که از آن
تغذیه کند بی‌اعتبار می‌سازد، و چارچوب اعتبارسنجیِ سهل‌گیر، هر عدد دقتی را که تولید
کند خوش‌بینانه می‌کند.

\cref{sec:3-2} شش فرض بنیادی پژوهش را \lr{—} همان‌گونه که در پیشنهاده‌ی مصوب صورت‌بندی
شده‌اند \lr{—} می‌آورد و از همان آغاز روشن می‌کند که شواهد فصل‌های چهارم و پنجم کدام
یک از آن‌ها را تأیید نکرده‌اند. \cref{sec:3-3} داده‌ها را معرفی می‌کند: مجموعه‌داده بنچمارک
\lr{EPEX-DE} که تمام نتایج پایان‌نامه بر آن استوار است، مجموعه‌ی داده‌ی زنده که تنها
برای آزمون برون‌توزیعی و ابزار عملیاتی به کار می‌رود، و تحلیل اکتشافی‌ای که
تصمیم‌های روش‌شناختیِ پس از خود را توجیه می‌کند. \cref{sec:3-4} مهندسی ویژگی و قید علّیِ
حاکم بر آن را تشریح می‌کند.

\cref{sec:3-5} چارچوب اعتبارسنجی، معیارهای دقت و آزمون معناداری را تعریف می‌کند. اینکه این
بخش پیش از معرفی مدل‌ها می‌آید تصادفی نیست: معیارها و پروتکل ارزیابی پیش از دیده‌شدن
هر نتیجه‌ای تثبیت شده‌اند، و همین پیش‌تعهد است که به اعداد فصل چهارم اعتبار می‌بخشد.
بخش‌های ۳-۶ و ۳-۷ به‌ترتیب مدل‌های پایه (ساده، \lr{SARIMAX}، \lr{LEAR-LASSO}) و
مدل‌های یادگیری ماشین و یادگیری عمیق (\lr{LightGBM}، \lr{LSTM}) را معرفی می‌کنند، و
\cref{sec:3-8} مدل ترکیبی و وزن‌دهیِ گروهیِ رژیم‌محور را.

سه قید در سرتاسر این فصل تکرار می‌شود و بهتر است از همین‌جا صریح باشد. نخست،
هم‌سانی: همه‌ی مدل‌ها از یک منبع یکتای ویژگی تغذیه می‌شوند، بر یک پروتکل یکسان رانده
می‌شوند و با یک پیمانه‌ی معیار سنجیده می‌شوند، تا اختلاف عملکرد به خودِ مدل بازگردد و
نه به تفاوت در شرایط. دوم، هم‌ارزی با پروتکل منتشرشده‌ی لاگو و همکاران
\cite{lago_forecasting_2021}، که مقایسه‌ی فصل چهارم با اعداد مرجع را ممکن می‌کند.
سوم، ثبت انحراف‌ها: هرجا پیاده‌سازی از پروتکل مرجع فاصله گرفته است \lr{—} به‌ویژه در
آهنگ واسنجیِ \lr{SARIMAX} و \lr{LSTM} \lr{—} این فاصله در متن تصریح شده است، به‌جای
آنکه در کد پنهان بماند.
